\section{Finite reduction and a linear cutoff}
\label{sec:bounds}

The repeated-exponent decomposition leaves two possible obstructions:
a nonsquare quadratic discriminant and a noninteger root of the symmetric
cubic. Exact factor pairs reduce the first exceptional set, while sign
intervals settle the large-exponent tail.

Fix $a,b\geq1$, and put
\[
h=a+1,\quad C=(a-1)(2a-1),\quad M=a^3(2a+1),\quad
A=h^2b+a(3a+1).
\]
The characteristic discriminant of~\eqref{eq:antisymmetric} is
\begin{equation}
\label{eq:discriminant}
D=h^2b^2+2a(3a+1)b+a^2.
\end{equation}
If $D$ is not a square, this block has irrational eigenvalues, which
remain irrational after the integer shift $a^2b-1$.

\begin{proposition}[Exact divisor reduction]
\label{prop:divisors}
For fixed $a\geq1$, the positive integers $b$ for which $D$ is a square
are exactly those obtained from positive factor pairs $d\leq e$ of $M$
such that
\begin{equation}
\label{eq:divisor-conditions}
b=\frac{d+e-a(3a+1)}{h^2}\in\mathbb Z_{\geq1},\qquad
s=\frac{e-d}{h}\in\mathbb Z_{\geq0}.
\end{equation}
In that case $D=s^2$.
\end{proposition}
\begin{proof}
The polynomial identity
\begin{equation}
\label{eq:product}
A^2-h^2D=4M
\end{equation}
holds for arbitrary $a,b$. If $D=s^2$ with $s\geq0$, then
$A-hs$ and $A+hs$ are positive: their product is positive and their sum
is $2A>0$. They have the same parity; since their product is divisible
by $4$, both are even. Write $A-hs=2d$ and $A+hs=2e$.
Then $de=M$, $d+e=A$ and $e-d=hs$, giving~\eqref{eq:divisor-conditions}.
Conversely those conditions yield $A=d+e$ and $hs=e-d$.
Using $(d+e)^2-(e-d)^2=4de=4M$ in~\eqref{eq:product} gives $D=s^2$.
\end{proof}

The divisor set is finite for every fixed $a$; this criterion has no
truncation in $b$. It gives first a quadratic bound.

\begin{corollary}
\label{cor:quadratic}
If $D$ is a square then $b\leq C$. Consequently every $b>C$ gives
a nonintegral Laplacian spectrum. For $a=1$, every $b\geq1$ is covered.
\end{corollary}
\begin{proof}
Since $(d-1)(e-1)\geq0$, $A=d+e\leq M+1$. The identity
\[
M+1-a(3a+1)=h^2C
\]
then gives $b\leq C$. The nonsquare case is covered by the antisymmetric
block. At $a=1$, $C=0$.
\end{proof}

\begin{lemma}
\label{lem:congruence}
Every admissible factor pair in Proposition~\ref{prop:divisors} satisfies
$d\equiv e\equiv1\pmod h$.
\end{lemma}
\begin{proof}
From $e-d=hs$ and $d+e=A$, reduction modulo $h$ gives
$e\equiv d$ and $2d\equiv2$, since $a(3a+1)\equiv2\pmod h$.
Also
\begin{equation}
\label{eq:strong-product}
(d-1)(e-1)=M-A+1=h^2(C-b).
\end{equation}
If $h$ is odd, the congruence $2(d-1)\equiv0$ gives the assertion.
If $h$ is even, the alternative to $d\equiv1$ is
$d-1\equiv h/2\pmod h$. Then $d-1$ and $e-1$ would both equal
$h/2$ times an odd integer. Their product divided by $h^2$ would be an
odd integer divided by $4$, contradicting~\eqref{eq:strong-product}.
This excludes the alternative for either parity of $a$.
\end{proof}

\begin{theorem}[Unified linear cutoff]
\label{thm:linear}
Let $a\geq2$ and
\[
R(a)=\frac{2(a-1)(a-2)}{a+2}=2a-10+\frac{24}{a+2}.
\]
For every integer $b>R(a)$ and distinct primes $p,q,r$, the graph
$G_{p^a q^a r^b}$ is not Laplacian integral. Every unresolved
repeated-exponent pair therefore has square $D$ and
$1\leq b\leq\lfloor R(a)\rfloor$.
\end{theorem}
\begin{proof}
If $D$ is nonsquare the antisymmetric block suffices. Otherwise use
Proposition~\ref{prop:divisors}. When $d=1$, $e=M$ and $b=C$;
Theorem~\ref{thm:boundary} settles this case.
If $d>1$, Lemma~\ref{lem:congruence} gives $d\geq T=a+2$.
Since $e\geq d\geq T$,
\[
\frac MT+T-(d+e)=(d-T)\left(\frac eT-1\right)\geq0.
\]
Hence
\[
b\leq\frac{M/T+T-a(3a+1)}{h^2}
=\frac{2(a-1)(a-2)}{a+2}=R(a).
\]
Thus $b>R(a)$ leaves either a nonsquare discriminant or the already
settled boundary case.
\end{proof}

At $(a,b)=(4,2)$ the square-discriminant pair has $d=6=a+2$,
so equality in the factor estimate is attained. The graph is still
nonintegral, as shown in the next section.

\subsection{Settling the first nonboundary factor}

The first possible value after $d=1$ is $d=a+2$.
Since $d$ divides $M$, the identity
\[
M=(a+2)(2a^3-3a^2+6a-12)+24
\]
forces $a+2\mid24$. With $a\geq2$ this leaves
$a=2,4,6,10,22$. Substitution into~\eqref{eq:divisor-conditions}
gives $b=0,2,5,12,35$, respectively. Excluding $b=0$, the complete
positive-$b$ list and its cubic certificates are:
\begin{center}
\begin{tabular}{clc}
\toprule
$(a,b)$ & $f_{a,b}(x)$ & Root-free prime\\
\midrule
$(4,2)$ & $x^3-86x^2+2304x-18816$ & 5\\
$(6,5)$ & $x^3-245x^2+19266x-488160$ & 7\\
$(10,12)$ & $x^3-842x^2+230160x-20534400$ & 13\\
$(22,35)$ & $x^3-4919x^2+7887858x-4132479120$ & 17\\
\bottomrule
\end{tabular}
\end{center}
Each polynomial has no root modulo its indicated prime. Its monic cubic
is therefore irreducible over $\mathbb Q$, and its real nonzero roots
lift to noninteger graph eigenvalues. Thus every possible $d=a+2$ pair
is settled, over all $a$, by this finite divisor argument.

\begin{theorem}[Second linear cutoff]
\label{thm:second}
For $a\geq2$ put
\[
S(a)=\frac{(a-3)(2a-3)}{2a+3}=a-6+\frac{27}{2a+3}.
\]
For every positive integer $b>S(a)$ and distinct primes $p,q,r$,
$G_{p^a q^a r^b}$ is not Laplacian integral.
Every remaining square-discriminant case has $d\geq2a+3$ and
$b\leq\lfloor S(a)\rfloor$.
\end{theorem}
\begin{proof}
Nonsquare $D$ is covered by the antisymmetric block. If $D$ is square,
the cases $d=1$ and $d=a+2$ are already settled. Every other admissible
factor satisfies $d\geq T=2a+3$ by Lemma~\ref{lem:congruence}.
Since $e\geq d\geq T$, the same factor inequality gives
\[
b\leq\frac{M/T+T-a(3a+1)}{(a+1)^2}
=\frac{(a-3)(2a-3)}{2a+3}=S(a).
\]
This proves the claim for $b>S(a)$.
\end{proof}

The new bound is strictly smaller than $R(a)$ for all $a\geq2$, since
\[
R(a)-S(a)=\frac{2a^3-a^2-a-6}{(a+2)(2a+3)}>0.
\]
Writing $z=a-2\geq0$, its numerator is $2z^3+11z^2+19z+4$.

\subsection{A finite reduction for each factor index}

For a nonboundary square-discriminant pair, write $d=1+hj$, $j\geq1$.
The following criterion fixes $j$ instead of $a$.

\begin{proposition}
\label{prop:index}
Fix $j\geq1$ and set $K_j=(j+1)^3(j+2)$.
Every admissible square-discriminant pair of index $j$ is obtained from
a positive divisor $d$ of $K_j$ with
$a=(d-1)/j-1\in\mathbb Z_{\geq2}$.
Set $e=M/d$, $v=(e-1)/(a+1)$, and retain exactly those candidates with
$d\leq e$ and $b=C-jv\geq1$. Then $e,v,b$ are integers and
$D=(v-j)^2$. Conversely every retained candidate is admissible.
\end{proposition}
\begin{proof}
Modulo $d=j(a+1)+1$ we have $ja\equiv-(j+1)$, whence
\[
j^4M=(ja)^3j(2a+1)\equiv(j+1)^3(j+2)=K_j\pmod d.
\]
Since $\gcd(d,j)=1$, we have $d\mid M$ if and only if $d\mid K_j$.
Thus $e$ is integral. Since $M\equiv d\equiv1\pmod h$, also
$e\equiv1\pmod h$, so $v$ is integral. The identity
\[
M=1+h(3a-2)+h^2C=(1+hj)(1+hv)
\]
gives $j+v+h jv=3a-2+hC$. Defining $b=C-jv$, we obtain
$d+e=h^2b+a(3a+1)$ and $e-d=h(v-j)$.
The ordering condition gives $v-j\geq0$; Proposition~\ref{prop:divisors}
now proves admissibility and $D=(v-j)^2$. The reverse direction follows
by applying these identities to any admissible index-$j$ pair.
\end{proof}

This is a complete finite problem for each fixed $j$, without a
truncation in $a$. Section~\ref{sec:completion} combines the first
three such problems with a uniform argument for all later indices.
For $j=2$, $K_2=108$ and $d=2a+3$ is an odd divisor of $108$ with
$d\geq7$. The only candidates are $d=9,27$, giving $a=3,12$.
The first has $e=21,v=5,b=0$ and is excluded. The second has
$e=1600,v=123,b=7$ and $\sqrt D=121$.
Thus $(12,7)$ is the only positive-$b$ case with $j=2$.

Its symmetric cubic is $f_{12,7}(x)=x^3-751x^2+180348x-13829760$.
Its values at all residues modulo $13$ are
$4,7,9,3,8,4,10,6,11,5,7,10,7$.
As a monic cubic with no root modulo $13$, it is irreducible over
$\mathbb Q$. Its real nonzero roots lift to noninteger graph eigenvalues,
settling the entire index-$2$ case.

\begin{theorem}[Third linear cutoff]
\label{thm:third}
For every $a\geq2$, put $T(a)=2(a-4)(a-2)/(3a+4)$.
If $b\geq1$ and $b>T(a)$, then $G_{p^a q^a r^b}$ is not Laplacian
integral. Every unresolved square-discriminant pair has $j\geq3$,
$d\geq3a+4$, and $b\leq\lfloor T(a)\rfloor$.
\end{theorem}
\begin{proof}
Nonsquare $D$ gives irrational antisymmetric eigenvalues.
For square $D$, the cases $j=0,1,2$ are all settled above.
Every other pair has $d\geq Q=3a+4$, and $e\geq d\geq Q$ gives
\[
b\leq\frac{M/Q+Q-a(3a+1)}{(a+1)^2}
=\frac{2(a-4)(a-2)}{3a+4}=T(a).
\]
This proves the claimed implication.
\end{proof}

For $a\geq3$ the new bound is strictly below $S(a)$, since
\[
S(a)-T(a)=\frac{2a^3-a^2-5a-12}{(2a+3)(3a+4)}>0.
\]
With $z=a-3\geq0$, the numerator is $2z^3+17z^2+43z+18$.
The exact expression
$T(a)=2a/3-44/9+320/(9(3a+4))$ shows its leading term.
